\section{Tujuan Penelitian}

Secara umum, penelitian ini bertujuan untuk merancang dan mengevaluasi CoDA-EDA (\textit{Compact Dependency-Aware Estimation of Distribution Algorithm}), sebuah algoritma metaheuristik yang ringan, sadar terhadap korelasi antar-atribut, dan mampu diperbarui secara inkremental, untuk mendukung \textit{policy mining Attribute-Based Access Control} (ABAC) pada perangkat \textit{edge}-IoT berdaya rendah.
Secara khusus, tujuan penelitian ini adalah untuk:
\begin{enumerate}
    \item Merancang representasi model probabilistik dan mekanisme pembelajaran dependensi atribut yang ringan namun mampu menangkap korelasi antar-atribut ABAC.
    \item Mengukur dan membandingkan kualitas kebijakan ABAC yang dihasilkan oleh algoritma yang diusulkan terhadap algoritma metaheuristik \textit{population-based} yang ada serta algoritma \textit{compact} murni pada dataset ABAC dengan tingkat korelasi antar-atribut yang bervariasi.
    \item Merancang dan menguji strategi pembaruan struktur dependensi secara inkremental yang dipicu oleh perubahan signifikan pada log akses, serta mengukur efisiensi komputasinya dibandingkan pendekatan pembaruan struktur konvensional (BOA/iBOA).
    \item Mengukur jejak memori dan waktu komputasi algoritma yang diusulkan ketika dijalankan pada perangkat \textit{edge}-IoT berdaya rendah, serta membandingkannya dengan algoritma metaheuristik ABAC \textit{population-based} yang ada saat ini.
\end{enumerate}
